1.7. If $Y$ is a complete intersection of dimension at least three in a
projective space $\mathbb P^r(k)$ over a field $k$, then
$\operatorname{Pic}(Y)$ is known to be isomorphic to $\mathbb Z$ and to be
generated by $\underline O(1)$ (SGA 2, XII, 3.7).  In dimension two, when
$Y$ is smooth, one can still prove the following.

\underline{Theorem} 1.8. \underline{Let $Y$ be a smooth complete-intersection
surface in the projective space $\mathbb P^r(k)$ over an algebraically
closed field $k$.  Then $\underline{\operatorname{Pic}}^\tau(Y)=0$, and the
class of $\underline O(1)$ is not divisible in the Néron--Severi group
$\operatorname{NS}(Y)$.}

Since $H^1(Y,\underline O)=0$ by 1.5(iii)$_a$, one knows that
$\underline{\operatorname{Pic}}^0(Y)=0$.  The vanishing of
$\underline{\operatorname{Pic}}^\tau(Y)$ then follows from the familiar
lemma below, the fact that $H^2(Y,\mathbb Z_\ell(1))$ is torsion-free by
1.6(iv), and the fact that $H^0(Y,\Omega_Y^1)=0$ by 1.5(iii)$_a$.

\underline{Lemma} 1.9. \underline{Let $Y$ be a complete algebraic variety
over an algebraically closed field $k$.}

\begin{enumerate}[label=(\roman*),leftmargin=2.6em]
\item \underline{If $\ell$ is prime to the characteristic of $k$, then the
$\ell$-torsion of $\operatorname{NS}(Y)$ identifies with the $\ell$-torsion
of $H^2(Y,\mathbb Z_\ell(1))$.}

%%%%%%%%%% source scan idx 57 | folio 50 %%%%%%%%%%

\item \underline{If $k$ has nonzero characteristic $p$ and $Y$ is smooth,
then the kernel of multiplication by $p$ in $\operatorname{Pic}(Y)$
identifies with the subgroup of $H^0(Y,\Omega_Y^1)$ consisting of locally
logarithmic differentials.}
\end{enumerate}

The Kummer exact sequences

\begin{equation*}
0\longrightarrow\mu_{\ell^n}\longrightarrow\mathbb G_m
\xrightarrow{\ x\mapsto x^{\ell^n}\ }\mathbb G_m\longrightarrow0
\end{equation*}

give exact sequences

\begin{equation*}
0\longrightarrow
\operatorname{Pic}(Y)/\ell^n\operatorname{Pic}(Y)
\longrightarrow H^2(Y,\mu_{\ell^n})
\longrightarrow H^2(Y,\mathbb G_m)_{\ell^n}
\longrightarrow0,
\end{equation*}

where

\begin{equation*}
\operatorname{Pic}(Y)/\ell^n\operatorname{Pic}(Y)
=\operatorname{NS}(Y)\otimes\mathbb Z/\ell^n.
\end{equation*}

Passing to the inverse limit gives

\begin{equation*}
0\longrightarrow\operatorname{NS}(Y)\otimes\mathbb Z_\ell
\longrightarrow H^2(Y,\mathbb Z_\ell(1))
\longrightarrow
\operatorname{Hom}\bigl(\mathbb Q_\ell/\mathbb Z_\ell,
H^2(Y,\mathbb G_m)\bigr)
\longrightarrow0.
\end{equation*}

The cokernel is torsion-free because $\mathbb Q_\ell/\mathbb Z_\ell$ is
divisible.  This proves (i).

To prove (ii), it is enough to observe that the exact sequence

\begin{equation*}
0\longrightarrow\underline O_Y^*
\xrightarrow{\ x\mapsto x^p\ }\underline O_Y^*
\xrightarrow{\ df/f\ }\Omega_Y^1
\end{equation*}

defines a short exact sequence

\begin{equation*}
0\longrightarrow\underline O_Y^*\longrightarrow\underline O_Y^*
\longrightarrow\Omega^1_{Y,\mathrm{loc}\,\log}\longrightarrow0,
\end{equation*}

and then to take the corresponding long exact cohomology sequence.  This
argument is due to Cartier, who proves moreover that, if $C$ denotes the
Cartier operator on $H^0(Y,\Omega_Y^1)$, then the locally logarithmic
differentials form the kernel of $C-\operatorname{Id}$.

Suppose, for a contradiction, that the class of $\underline O(1)$ in
$\operatorname{NS}(Y)$ has the form $nu$, with $n>1$.  If $\ell$ divides
$n$ and is prime to the characteristic, then in $\ell$-adic cohomology one
again has

%%%%%%%%%% source scan idx 58 | folio 51 %%%%%%%%%%

\begin{equation*}
\eta_\ell=n\,c_1(u),
\end{equation*}

contradicting 1.6(iv).

Finally, if the characteristic $p$ of $k$ is nonzero and $p$ divides $n$,
then in Hodge cohomology one again has

\begin{equation*}
\eta=n\,c_1(u)=0,
\end{equation*}

contradicting 1.5(iii)$_d$ and completing the proof.

